% --- ARCHIVO: Capitulos/introduccion.tex ---

\chapter{Bienvenidos a la Era del Martillo}

\begin{loreblock}
``En un mundo donde la energía de la Gema se agota, cada paso sobre la tierra es un robo al destino. Los Biskrorums marchan con acero, los Evolgens con voluntad, pero todos buscan lo mismo: el control del flujo. Cuando el martillo cae, no solo se rompen huesos — se quiebra la realidad.''
\end{loreblock}

\vspace{0.4cm}

\section{¿Qué es La Era del Martillo?}

\textbf{La Era del Martillo} es un juego de escaramuzas heroicas diseñado para dos jugadores que buscan profundidad táctica y consecuencias reales en cada decisión. A diferencia de otros juegos de guerra, aquí cada unidad es un activo invaluable. El sistema no se trata solo de mover miniaturas, sino de gestionar el agotamiento, la posición y, sobre todo, la energía del entorno.

\begin{reglaespecial}{La Filosofía del Juego}
Este sistema se basa en tres pilares fundamentales:
\begin{itemize}
    \item \textbf{Letalidad Detallada:} Gracias al sistema de localización de daños, un solo golpe bien colocado puede cambiar el curso de la batalla. Los miembros inutilizados, las heridas en la cabeza, todo tiene consecuencias permanentes.
    \item \textbf{Gestión de Recursos:} Los Puntos de Acción (PA) y Energía (PE) dictan el ritmo. Un guerrero agotado es un guerrero muerto.
    \item \textbf{Zonas Secas:} El uso irresponsable de la magia tiene un costo físico en el campo de batalla, alterando el mapa de forma permanente e irreversible.
\end{itemize}
\end{reglaespecial}

\section{Lo que Necesitas para Jugar}

\begin{description}
    \item[Miniaturas:] Que representen a tus héroes y escuadras. Cada unidad tiene su ficha correspondiente en el Codex de tu facción.
    \item[Dados:] Un dado de doce caras (\textbf{1d12}) para pruebas de éxito, un dado de cien (\textbf{1d100}) para localización, un dado de seis (\textbf{1d6}) para efecto de estado, y ocasionalmente \textbf{1d4}.
    \item[Superficie de Juego:] Un área de \textbf{90cm $\times$ 90cm} con escenografía que represente ruinas, bosques o zonas de extracción de energía.
    \item[Cinta Métrica:] Las distancias se miden en centímetros (cm).
    \item[Fichas de Turno y Marcadores:] Para registrar el estado Agotado, Zonas Secas, efectos de estado y puntos de vida.
    \item[Mazos de Cartas:] Un mazo de \textbf{Misiones Principales} (cartas grandes, formato tipo Magic) y un mazo de \textbf{Misiones Secundarias} (secretas).
    \item[Cartas de Lista:] La lista de tu ejército de hasta 60 puntos.
\end{description}

\begin{imagebox}{Lo que necesitas para jugar — dados, fichas y miniaturas sobre la mesa de batalla}
    \vspace{2cm}
\end{imagebox}

\section{El Concepto de Escaramuza Heroica}

A diferencia de las grandes batallas de bloques de infantería, este manual se enfoca en el \textbf{Combate Heroico de Escaramuzas}. Cada miniatura actúa de forma independiente y cuenta con acciones, reacciones y habilidades únicas. Un solo Héroe puede cambiar el resultado de una batalla — pero también puede caer ante una táctica bien ejecutada.

Los cuerpos de las unidades caídas \textbf{permanecen en el tablero}, ya que pueden interferir en el movimiento, ser usados como cobertura, o incluso interactuar con ciertas habilidades de facción.

\begin{tcolorbox}[colback=HammerRed!8, colframe=HammerRed, title=\textbf{NOTA DEL DISEÑADOR}]
\small Este juego está pensado para partidas de 60--90 minutos una vez que ambos jugadores conocen sus facciones. La primera partida puede extenderse mientras se aprenden las mecánicas. ¡Se recomienda una partida de aprendizaje con listas prearmadas!
\end{tcolorbox}
